% ============================================================
% CAS/Python ενότητες — Κεφάλαιο 11: Ορισμένο Ολοκλήρωμα
%
% QR code: qr_b29000ea8782.png  (→ latex/qrcodes/)
% URL: https://nikmatz.github.io/ELADIC_GR/code/ch11/
% ============================================================

\casactivbox%
  {Maxima: Ορισμένο Ολοκλήρωμα — ΘΘΛ, Εμβαδόν, Αριθμητική Ολοκλήρωση}%
  {%
    \label{cas:ch11_defintegral}
    Χρησιμοποιήστε \texttt{integrate(f,x,a,b)} για τα παρακάτω.
    \begin{enumerate}[label=(\alph*),itemsep=2pt]
      \item \textbf{Θεμελιώδες Θεώρημα Λογισμού}: υπολογίστε
        $\int_0^3 x^2\,dx$, $\int_0^{\pi}\sin x\,dx$,
        $\int_0^1 e^x\,dx$, $\int_1^e \frac{1}{x}\,dx$.
        Επαληθεύστε μέσω $F(b)-F(a)$ χρησιμοποιώντας
        \texttt{subst} και \texttt{float}.
      \item \textbf{Ιδιότητες}: επαληθεύστε γραμμικότητα,
        αντιστροφή ορίων $\int_a^b = -\int_b^a$
        και πρόσθεση $\int_a^b = \int_a^c + \int_c^b$.
      \item \textbf{Εμβαδόν μεταξύ καμπυλών}:
        βρείτε τομές με \texttt{solve}, υπολογίστε
        $A = \int_{\alpha}^{\beta}[f_{\text{άνω}}(x)-f_{\text{κάτω}}(x)]\,dx$
        για $y=x$ vs $y=x^2$ και $y=-x^2+4$ vs $y=x^2-2x$.
      \item \textbf{Αριθμητική ολοκλήρωση}: χρησιμοποιήστε
        \texttt{romberg} για $\int_0^1 e^{x^2}\,dx$ και
        $\int_0^1\sin(x^2)\,dx$.
        Συγκρίνετε με κανόνα τραπεζίου $n=4$.
    \end{enumerate}
    \smallskip
    \textbf{Εντολές Maxima:}
    \texttt{integrate(f,x,a,b)}, \texttt{romberg(f,x,a,b)},
    \texttt{quad\_qags(f,x,a,b)}, \texttt{solve(f=g,x)},
    \texttt{float()}
  }%
  {https://nikmatz.github.io/ELADIC_GR/code/ch11/}%
  {qr_b29000ea8782.png}

\subsection*{Δραστηριότητα Python}

\begin{pybox}{Python: Ορισμένο Ολοκλήρωμα — SymPy \& SciPy}%
             {https://nikmatz.github.io/ELADIC_GR/code/ch11/}%
             {qr_b29000ea8782.png}
\label{py:ch11_defintegral}
\begin{enumerate}[label=(\alph*),itemsep=2pt]
  \item Υπολογίστε με \texttt{sympy.integrate(f,(x,a,b))} τα
    $\int_0^3 x^2\,dx$, $\int_0^{\pi}\sin x\,dx$,
    $\int_0^1 e^x\,dx$, $\int_1^e\frac{1}{x}\,dx$,
    $\int_0^1\frac{1}{x^2+1}\,dx$.
    Συγκρίνετε ακριβείς τιμές με \texttt{float()}.
  \item Επαληθεύστε τις ιδιότητες (γραμμικότητα, πρόσθεση
    διαστημάτων, αντιστροφή ορίων) αριθμητικά.
  \item Εμβαδόν μεταξύ καμπυλών: για $y=x$ vs $y=x^2$
    σχεδιάστε τις καμπύλες, σκιάστε την ενδιάμεση περιοχή
    (\texttt{fill\_between}) και εκτυπώστε το εμβαδόν.
  \item Αριθμητική ολοκλήρωση με \texttt{scipy.integrate.quad}
    για υπερβατικά ολοκληρώματα ($e^{x^2}$, $\sin(x^2)$).
    Συγκρίνετε σφάλματα κανόνα τραπεζίου vs Simpson
    για $n = 4, 8, 16, 32$ σε γράφημα σύγκλισης (log-log).
\end{enumerate}
\medskip
\textbf{Βιβλιοθήκες / Εντολές:}
\texttt{sympy.integrate()}, \texttt{scipy.integrate.quad()},
\texttt{np.trapz()}, \texttt{scipy.integrate.simpson()},
\texttt{matplotlib.fill\_between()}
\end{pybox}
